% Part: turing-machines
% Chapter: machines-computations
% Section: church-turing-thesis

\documentclass[../../../include/open-logic-section]{subfiles}

\begin{document}

\olfileid{tur}{mac}{ctt}
\olsection{أطروحة تشرتش--تورنغ}

المقصود بآلة تورنغ أن تضبط معنى الإجراء الفعال ضبطًا رياضيًّا. وقد رأى تورنغ أن من استوعب المعنيين يدرك حدسيًّا أن كل ما ينجزه إجراء فعال تنجزه آلة تورنغ. ومما يقوي هذه الدعوى أن سائر البدائل الدقيقة المقترحة للإجراء الفعال انتهت إلى التكافؤ الامتدادي مع آلات تورنغ، أي إلى حساب الدوال نفسها من غير زيادة ولا نقصان. وتعرف هذه الدعوى بـ\emph{أطروحة تشرتش--تورنغ}.

\begin{defn}[أطروحة تشرتش--تورنغ]
\emph{أطروحة تشرتش--تورنغ}: كل ما يحسب بإجراء فعال قابل للحساب بتورنغ.
\end{defn}

وللاحتجاج بالأطروحة وجهان. أولهما إعفاء النفس من عناء البناء التفصيلي: إذا وصفنا إجراء فعالًا يحسب شيئًا، بعبارة تشبه «الشفرة الكاذبة» مثلًا، احتججنا بالأطروحة لوجود آلة تورنغ تحسب الدالة نفسها، وإن لم ننشئ الآلة بالفعل.

والثاني أعمق دلالة فلسفية. فقد نبرهن أن دالة لا تحسبها أي آلة تورنغ، ثم نستعين بالأطروحة لنستنتج امتناع حسابها بأي إجراء فعال، كائنًا ما كان. إذ لو وجد الإجراء لوجدت له، بمقتضى الأطروحة، آلة تورنغ؛ وقد نفينا وجودها بالبرهان. ومن هنا يستدل بها خاصة على أن مسألة التوقف لا تعجز عنها آلات تورنغ وحدها، بل لا سبيل فعال إلى حلها أصلًا.

\end{document}
